% GENERATED FILE. Edit canonical lessons, then run npm run generate:books.
% canonical-source-hash: fnv1a64:a496c25f2f2f641f
% canonical-lessons: ZH-C47-nian, ZH-W47-nian, ZH-C47-yue, ZH-C47-mingnian, ZH-W47-ming-bright, ZH-C47-jinnian, ZH-C47-riqi, ZH-R47-review-43-47

\chapter{Months and Years}
\label{ch:zh-months-and-years}
\hlchaptermodality{47}

\begin{chapteropening}
\textbf{By the end of this chapter:} \emph{I can say this year, next year, the month and the date.}

Complete the last lesson of chapter 47: I can say this year, next year, the month and the date.
\end{chapteropening}

\section[\texorpdfstring{\zh{年} (nián)}{年 (nián)}]{\zh{年} (nián) — year}

\label{lesson:ZH-C47-nian}



\begin{quote}
Before the new one: say the Mandarin for weekend, then the Mandarin for Sunday.
\end{quote}

\subsection*{You'll want to know: \zh{年}}
\textbf{\zh{年}} — \emph{nián} — \textquotedblleft{}year\textquotedblright{}.

\begin{quote}
\textbf{\zh{一年有多少个星期}} \emph{yì nián yǒu duōshao ge xīngqī} — How many weeks are there in a year?
\end{quote}

The new character, in pieces: \zh{年} — six strokes: a short fall, two bars, a short vertical on the left, a long bar, and one long vertical through them. The next lesson puts it in your hand.

\begin{cousinweb}
\textbf{\zh{年}} (\emph{nián}) is inherited from Middle Chinese \emph{nen} (Guangyun, in Baxter's transcription), glossed 'the grain is ripe'. One harvest a year is the usual explanation of the sense, though the step is not settled.
\end{cousinweb}

\begin{grammarlens}[title={What you've built}]
One more word for times and partings.
\end{grammarlens}

\subsection*{Your turn}
\begin{itemize}
\raggedright
  \item \emph{Say it:} \emph{nián}
  \item \emph{Say it:} \emph{nián}, once more
  \item \emph{Say it:} \emph{yì nián yǒu duōshao ge xīngqī} — How many weeks are there in a year?
  \item \emph{Recall:} say the Mandarin for weekend, then the Mandarin for Sunday, then say \emph{nián} again
\end{itemize}

\begin{tcolorbox}[breakable,colback=teal!4,colframe=teal!35!black,title={Before you move on}]
What does \zh{年} mean? (\textquotedblleft{}Year\textquotedblright{}.) Say it once more.
\end{tcolorbox}
\section[\texorpdfstring{\zh{年} (nián)}{年 (nián)}]{\zh{年} — three bars and a long vertical}

\label{lesson:ZH-W47-nian}



\begin{quote}
The word first, then the character inside it.

\begin{itemize}
\raggedright
  \item \emph{Recall:} say \emph{nián} — \textbf{R1}, one lesson back, before it decays
\end{itemize}

\begin{quote}
\zh{生} \zh{十}
\end{quote}

\emph{Read it:} each one aloud, with nothing said first — \emph{shēng}, \emph{shí}
\end{quote}

\subsection*{Script — three bars and a long vertical}
\hlblockfigure{\detokenize{figures/ZH-W47-nian-filmstrip.pdf}}{How \zh{年} is written, stroke by stroke}

\begin{quote}
\zh{年}
\end{quote}

\textbf{One piece, six strokes.}

\textbf{A short left-falling stroke} at the top left.

\textbf{Two horizontals}, the upper one short, the one below it longer.

\textbf{A short vertical} down from the left end of the second bar.

\textbf{A long horizontal} across the bottom.

\textbf{Last, the long vertical} down through the middle, past the bottom bar.

\subsection*{Your turn}
\begin{itemize}
\raggedright
  \item \emph{Write it:} \zh{年}, all six, with the model in front of you
  \item \emph{Cover:} write it again from memory, and count the pieces, not the strokes
\end{itemize}

\begin{tcolorbox}[breakable,colback=teal!4,colframe=teal!35!black,title={Before you move on}]
How many strokes in \zh{年}? (\textbf{Six}.) What comes first? (\textbf{A short left-falling stroke}.) Which stroke is last? (\textbf{The long vertical through the middle}.)

Source: \href{https://raw.githubusercontent.com/chanind/hanzi-writer-data/68d10a4b21150cae5e1ebbd223eed289cf32d90c/data/\%E5\%B9\%B4.json}{Hanzi Writer Data \zh{年}, pinned stroke paths and medians}.
\end{tcolorbox}
\section[yuè]{yuè — month; moon}

\label{lesson:ZH-C47-yue}



\begin{quote}
Before the new one: say the Mandarin for Sunday, then the Mandarin for year.
\end{quote}

\subsection*{You'll want to know: yuè}
\emph{yuè} — \textquotedblleft{}month; moon\textquotedblright{}.

\begin{quote}
\emph{zhège yuè wǒ hěn máng} — I'm busy this month.
\end{quote}

Its characters come later; for now it is a word for the ear.

\begin{cousinweb}
\emph{yuè} is inherited from Middle Chinese \emph{ngjwot} (Guangyun, in Baxter's transcription), where the dictionary quotes an old saying that the moon is the essence of water.
\end{cousinweb}

\begin{grammarlens}[title={What you've built}]
One more word for times and partings.
\end{grammarlens}

\subsection*{Your turn}
\begin{itemize}
\raggedright
  \item \emph{Say it:} \emph{yuè}
  \item \emph{Say it:} \emph{yuè}, once more
  \item \emph{Say it:} \emph{zhège yuè wǒ hěn máng} — I'm busy this month.
  \item \emph{Recall:} say the Mandarin for Sunday, then the Mandarin for year, then say \emph{yuè} again
\end{itemize}

\begin{tcolorbox}[breakable,colback=teal!4,colframe=teal!35!black,title={Before you move on}]
What does yuè mean? (\textquotedblleft{}Month; moon\textquotedblright{}.) Say it once more.
\end{tcolorbox}
\section[\texorpdfstring{\zh{明年} (míngnián)}{明年 (míngnián)}]{\zh{明年} (míngnián) — next year}

\label{lesson:ZH-C47-mingnian}



\begin{quote}
Before the new one: say the Mandarin for year, then the Mandarin for month.
\end{quote}

\subsection*{You'll want to know: \zh{明年}}
\textbf{\zh{明年}} — \emph{míngnián} — \textquotedblleft{}next year\textquotedblright{}.

\begin{quote}
\textbf{\zh{明年我买车}} \emph{míngnián wǒ mǎi chē} — I'm buying a car next year.
\end{quote}

The new character, in pieces: \zh{明} — \zh{日} on the left and, on the right, the moon piece, the word yuè from the lesson before. The next lesson puts it in your hand.

\begin{cousinweb}
\textbf{\zh{明年}} (\emph{míngnián}) is built from \zh{明} 'bright' and \zh{年} 'year', both inherited from Middle Chinese, \emph{mjaeng} and \emph{nen} (Guangyun, in Baxter's transcription); the dictionary glosses \emph{mjaeng} 'light'.
\end{cousinweb}

\begin{grammarlens}[title={What you've built}]
One more word for times and partings.
\end{grammarlens}

\subsection*{Your turn}
\begin{itemize}
\raggedright
  \item \emph{Say it:} \emph{míngnián}
  \item \emph{Say it:} \emph{míngnián}, once more
  \item \emph{Say it:} \emph{míngnián wǒ mǎi chē} — I'm buying a car next year.
  \item \emph{Recall:} say the Mandarin for year, then the Mandarin for month, then say \emph{míngnián} again
\end{itemize}

\begin{tcolorbox}[breakable,colback=teal!4,colframe=teal!35!black,title={Before you move on}]
What does \zh{明年} mean? (\textquotedblleft{}Next year\textquotedblright{}.) Say it once more.
\end{tcolorbox}
\section[\texorpdfstring{\zh{明} (míng)}{明 (míng)}]{\zh{明} — the sun beside the moon}

\label{lesson:ZH-W47-ming-bright}



\begin{quote}
The word first, then the character inside it.

\begin{itemize}
\raggedright
  \item \emph{Recall:} say \emph{míngnián} — \textbf{R1}, one lesson back, before it decays
\end{itemize}

\begin{quote}
\zh{日} \zh{期}
\end{quote}

\emph{Read it:} each one aloud, with nothing said first — \emph{rì}, \emph{qī}
\end{quote}

\subsection*{Script — the sun beside the moon}
\hlblockfigure{\detokenize{figures/ZH-W47-ming-bright-filmstrip.pdf}}{How \zh{明} is written, stroke by stroke}

\begin{quote}
\zh{明}
\end{quote}

\textbf{Two pieces, left before right: the sun, then the moon.}

\textbf{\zh{日}, four strokes}, narrow and set a little high on the left.

\textbf{The moon, four strokes}, as on the right of \zh{期}: a left-falling stroke, a horizontal that turns down and hooks, and two short horizontals inside.

4 + 4 = eight.

\subsection*{Your turn}
\begin{itemize}
\raggedright
  \item \emph{Write it:} \zh{明}, all eight, with the model in front of you
  \item \emph{Cover:} write it again from memory, and count the pieces, not the strokes
\end{itemize}

\begin{tcolorbox}[breakable,colback=teal!4,colframe=teal!35!black,title={Before you move on}]
Which two pieces make \zh{明}? (\textbf{\zh{日} and the moon}.) Which character from the last chapters has the same moon on the right? (\textbf{\zh{期}}.) How many strokes? (\textbf{Eight}.)

Source: \href{https://raw.githubusercontent.com/chanind/hanzi-writer-data/68d10a4b21150cae5e1ebbd223eed289cf32d90c/data/\%E6\%98\%8E.json}{Hanzi Writer Data \zh{明}, pinned stroke paths and medians}.
\end{tcolorbox}
\section[\texorpdfstring{\zh{今年} (jīnnián)}{今年 (jīnnián)}]{\zh{今年} (jīnnián) — this year}

\label{lesson:ZH-C47-jinnian}



\begin{quote}
Before the new one: say the Mandarin for month, then the Mandarin for next year.
\end{quote}

\subsection*{You'll want to know: \zh{今年}}
\textbf{\zh{今年}} — \emph{jīnnián} — \textquotedblleft{}this year\textquotedblright{}.

\begin{quote}
\textbf{\zh{我今年十八岁}} \emph{wǒ jīnnián shíbā suì} — I'm eighteen this year.
\end{quote}

\begin{cousinweb}
\textbf{\zh{今年}} (\emph{jīnnián}) is built from \zh{今} 'now' and \zh{年} 'year', both inherited from Middle Chinese, \emph{kim} and \emph{nen} (Guangyun, in Baxter's transcription).
\end{cousinweb}

\begin{grammarlens}[title={What you've built}]
One more word for times and partings.
\end{grammarlens}

\subsection*{Your turn}
\begin{itemize}
\raggedright
  \item \emph{Say it:} \emph{jīnnián}
  \item \emph{Say it:} \emph{jīnnián}, once more
  \item \emph{Say it:} \emph{wǒ jīnnián shíbā suì} — I'm eighteen this year.
  \item \emph{Recall:} say the Mandarin for month, then the Mandarin for next year, then say \emph{jīnnián} again
\end{itemize}

\begin{tcolorbox}[breakable,colback=teal!4,colframe=teal!35!black,title={Before you move on}]
What does \zh{今年} mean? (\textquotedblleft{}This year\textquotedblright{}.) Say it once more.
\end{tcolorbox}
\section[\texorpdfstring{\zh{日期} (rìqī)}{日期 (rìqī)}]{\zh{日期} (rìqī) — date}

\label{lesson:ZH-C47-riqi}



\begin{quote}
Before the new one: say the Mandarin for next year, then the Mandarin for this year.
\end{quote}

\subsection*{You'll want to know: \zh{日期}}
\textbf{\zh{日期}} — \emph{rìqī} — \textquotedblleft{}date\textquotedblright{}.

\begin{quote}
\textbf{\zh{请看日期}} \emph{qǐng kàn rìqī} — Please check the date.
\end{quote}

\begin{cousinweb}
\textbf{\zh{日期}} (\emph{rìqī}) is built from \zh{日} 'day' and \zh{期} 'appointed time', both inherited from Middle Chinese, \emph{nyit} and \emph{gi} (Guangyun, in Baxter's transcription).
\end{cousinweb}

\begin{grammarlens}[title={What you've built}]
One more word for times and partings.
\end{grammarlens}

\subsection*{Your turn}
\begin{itemize}
\raggedright
  \item \emph{Say it:} \emph{rìqī}
  \item \emph{Say it:} \emph{rìqī}, once more
  \item \emph{Say it:} \emph{qǐng kàn rìqī} — Please check the date.
  \item \emph{Recall:} say the Mandarin for next year, then the Mandarin for this year, then say \emph{rìqī} again
\end{itemize}

\begin{tcolorbox}[breakable,colback=teal!4,colframe=teal!35!black,title={Before you move on}]
What does \zh{日期} mean? (\textquotedblleft{}Date\textquotedblright{}.) Say it once more.
\end{tcolorbox}
\section[(dialogue)]{Looking back: from Coming and Going to Months and Years}

\label{lesson:ZH-R47-review-43-47}



\begin{quote}
Say the Mandarin for to go and to go back.
\end{quote}

\subsection*{Your turn}
\begin{itemize}
\raggedright
  \item \emph{Say it:} to go, to go back, to walk, to go out, school
  \item \emph{Say it:} car, restaurant, hospital, to sit, taxi
  \item \emph{Say it:} day, today, yesterday, now, tomorrow
  \item \emph{Say it:} week, Monday, moment, weekend, Sunday
  \item \emph{Say it:} year, month, next year, this year, date
  \item \emph{Read it:} \zh{去} \zh{车} \zh{天} \zh{今} \zh{星} \zh{期} \zh{年} \zh{明} — each one aloud, then a word that holds it
\end{itemize}

\begin{tcolorbox}[breakable,colback=teal!4,colframe=teal!35!black,title={Before you move on}]
To go? (\textbf{\zh{去}}, \emph{qù}.) To go back? (\emph{huí}.) To walk? (\emph{zǒu}.) To go out? (\emph{chūmén}.) School? (\emph{xuéxiào}.) Car? (\textbf{\zh{车}}, \emph{chē}.) Restaurant? (\emph{fàndiàn}.) Hospital? (\emph{yīyuàn}.) To sit? (\emph{zuò}.) Taxi? (\emph{chūzūchē}.) Day? (\textbf{\zh{天}}, \emph{tiān}.) Today? (\textbf{\zh{今天}}, \emph{jīntiān}.) Yesterday? (\emph{zuótiān}.) Now? (\emph{xiànzài}.) Tomorrow? (\emph{míngtiān}.) Week? (\textbf{\zh{星期}}, \emph{xīngqī}.) Monday? (\textbf{\zh{星期一}}, \emph{xīngqīyī}.) Moment? (\emph{shíhou}.) Weekend? (\emph{zhōumò}.) Sunday? (\textbf{\zh{星期天}}, \emph{xīngqītiān}.) Year? (\textbf{\zh{年}}, \emph{nián}.) Month? (\emph{yuè}.) Next year? (\textbf{\zh{明年}}, \emph{míngnián}.) This year? (\textbf{\zh{今年}}, \emph{jīnnián}.) Date? (\textbf{\zh{日期}}, \emph{rìqī}.)
\end{tcolorbox}
